\section{Data}

\subsection{Voteview DW-NOMINATE}

We use the Voteview.com DW-NOMINATE first-dimension ideal point estimates
for all House members from the 107th through 118th Congresses (2001--2024),
downloaded from \url{https://voteview.com/static/data/out/members/HSall\_members.csv}~\citep{lewis2024voteview}.
The dataset contains 5{,}355 member-congress records, excluding members
with fewer than 10 votes (insufficient data).

DW-NOMINATE dim1 ranges approximately from $-1$ (most liberal)
to $+1$ (most conservative).
We use the absolute value $|\text{dim1}|$ as the polarization measure
(higher = more extreme from the centre).

\subsection{Compactness Data}

District Polsby-Popper scores by congress are derived from the corresponding
redistricting cycle's enacted boundary files.
For the 107th--112th Congresses (2001--2012): the 2000 cycle enacted maps.
For the 113th--117th Congresses (2013--2022): the 2010 cycle enacted maps.
For the 118th Congress (2023--2024): the 2020 cycle enacted maps.
\textsc{bisect} compactness scores are from the 2020 algorithmic maps.

\subsection{District Partisan Composition}

We control for district partisan composition using the presidential vote
share from the most recent presidential election preceding each Congress.
Source: VEST precinct returns areal-interpolated to districts.
For states where full VEST data is unavailable, we use MIT Election Lab
county presidential returns.
